% L'impôt sur le revenu le plus efficace possible (note courte, complément de note-ir-csg.tex)
% Compilation depuis fiscalite/ir_csg/ : pdflatex note-efficace.tex (deux passes)
% Chiffres : python3 ir_csg/systeme_efficace.py -> ir_csg/sorties/systeme_efficace.txt (étiquettes [Sn])
% Figure : FIG_NU=1 python3 ir_csg/figures_efficace.py (depuis fiscalite/)
\documentclass[11pt,a4paper]{article}
% Préambule commun au brief et à la revue de littérature sur la taxe sur la valeur foncière.
% Même composition que le cours (cours/refaire-l-impot.tex) ; figures « nues » produites par
%   FIG_NU=1 python lvt/graphiques_lvt.py
\usepackage[T1]{fontenc}
\usepackage[utf8]{inputenc}
\usepackage{lmodern}
\usepackage[french]{babel}
\usepackage[a4paper,left=2.8cm,right=2.8cm,top=2.6cm,bottom=2.8cm]{geometry}
\usepackage{graphicx}
\usepackage[export]{adjustbox}
\usepackage{booktabs}
\usepackage{tabularx}
\usepackage{array}
\usepackage{xcolor}
\usepackage{titlesec}
\usepackage[labelfont=bf,labelsep=period,font=small,singlelinecheck=false,
            justification=raggedright,skip=6pt]{caption}
\usepackage{fancyhdr}
\usepackage{enumitem}
\usepackage{microtype}
\usepackage{textcomp}
\usepackage{amsmath}
\providecommand{\euro}{\texteuro}
\usepackage[hidelinks]{hyperref}

\definecolor{encre}{HTML}{1A1A19}
\definecolor{bleu}{HTML}{003399}
\definecolor{vert}{HTML}{1F7A4D}
\definecolor{ocre}{HTML}{C79100}
\definecolor{brique}{HTML}{993333}
\definecolor{filet}{HTML}{C8C6BE}

\titleformat{\section}{\normalfont\large\bfseries\color{bleu}}{\thesection.}{0.6em}{}
\titleformat{\subsection}{\normalfont\normalsize\bfseries\color{bleu}}{\thesubsection}{0.6em}{}
\titlespacing*{\section}{0pt}{2.4ex plus 1ex minus .2ex}{1.2ex plus .2ex}
\titlespacing*{\subsection}{0pt}{2ex plus 1ex minus .2ex}{0.8ex plus .2ex}

\pagestyle{fancy}
\fancyhf{}
\renewcommand{\headrulewidth}{0.4pt}
\fancyhead[R]{\small\thepage}

\setlength{\parindent}{0pt}
\setlength{\parskip}{0.62em}
\linespread{1.05}
\renewcommand{\arraystretch}{1.2}
\setlist{itemsep=2pt,topsep=2pt,parsep=0pt,leftmargin=1.4em}

\newenvironment{encadre}[1]
  {\par\medskip\noindent\begin{minipage}{\linewidth}%
   \color{encre}\setlength{\parskip}{0.5em}%
   \hrule height 0.8pt width \linewidth \vspace{6pt}%
   {\footnotesize\bfseries\color{vert}\MakeUppercase{#1}}\par\vspace{-2pt}%
   \small}
  {\par\vspace{4pt}\hrule height 0.4pt width \linewidth\end{minipage}\par\medskip}

\newcommand{\figbloc}[3]{%
  \begin{figure}[htbp]\centering
  \caption{#2}\label{#3}
  \includegraphics[max width=\linewidth,max height=0.40\textheight]{../figures/nu/#1.pdf}
  \IfFileExists{../figures/nu/#1_note.tex}{\par\vspace{4pt}%
    \parbox{\linewidth}{\raggedright\scriptsize\color{encre}\setlength{\parskip}{2pt}%
    \input{../figures/nu/#1_note.tex}}}{}%
  \end{figure}}

\newcommand{\pc}{\,\%}

\geometry{left=2.3cm,right=2.3cm,top=2.1cm,bottom=2.3cm}
\setlength{\parskip}{0.5em}
\renewcommand{\figbloc}[3]{%
  \begin{figure}[!htbp]\centering
  \caption{#2}\label{#3}
  \includegraphics[max width=\linewidth,max height=0.34\textheight]{../figures/nu/#1.pdf}
  \par\vspace{3pt}\parbox{\linewidth}{\raggedright\scriptsize\color{encre}\setlength{\parskip}{2pt}%
    \input{../figures/nu/#1_note.tex}}%
  \end{figure}}
\renewcommand{\topfraction}{0.9}\renewcommand{\textfraction}{0.08}\renewcommand{\floatpagefraction}{0.85}
\titlespacing*{\section}{0pt}{1.6ex plus .6ex minus .2ex}{0.8ex plus .2ex}
\fancyhead[L]{\small\itshape L'impôt sur le revenu le plus efficace possible}

\begin{document}
\renewcommand{\figurename}{Figure}\renewcommand{\tablename}{Tableau}

\thispagestyle{empty}
{\footnotesize\color{ocre}\bfseries NOTE DE TRAVAIL\quad·\quad NOVEMBRE 2026\par}
\vspace{6pt}
{\fontsize{20}{24}\selectfont\bfseries\color{bleu}L'impôt sur le revenu le plus efficace possible\par}
\vspace{5pt}
{\large Un taux unique sur le travail, un crédit par enfant, à recette inchangée\par}
\vspace{6pt}
{\small Colin Baget. Mêmes données et même modèle que la note « Fusionner les impôts sur le revenu »
(DGFiP, IRCOM et déclarations 2042, revenus 2024 ; OpenFisca-France pour les foyers types) ; programme
\texttt{ir\_csg/systeme\_efficace.py}.\par}
\vspace{5pt}
{\color{bleu}\hrule height 1.2pt}

\section*{Ce que l'on cherche}

La question est volontairement étroite : quel impôt unique, à la place de l'IR, de la CSG et de la CRDS
sur les revenus du travail et de remplacement, rapporte les mêmes 220~Md€ en décourageant le moins
possible l'activité, tout en aidant les familles ? L'égalité n'est pas un objectif ; la redistribution
n'est mesurée que pour information.

Un impôt sur le revenu coûte à l'économie plus que ce qu'il rapporte, parce que chacun peut réduire son
revenu imposable : travailler moins, refuser des heures ou une promotion, se faire payer autrement ou
plus tard. Ce qui pèse sur ces décisions est le taux marginal, la part prélevée sur le prochain euro, et
la perte qui en résulte (la perte sèche) croît à peu près comme son carré. D'où deux règles : un même
taux pour tous sur l'assiette la plus large possible (à recette égale, imposer deux moitiés de
l'assiette à 10 et 30~\% coûte un quart de plus qu'un taux unique de 20~\%, puisque $10^2+30^2=1\,000$
contre $2\times20^2=800$) ; des taux plus élevés sur ce qui ne réagit pas.

Poussée au bout, cette logique mène à un impôt par tête : environ 4~000~€ par adulte, sans aucun taux
marginal. Personne ne réduirait son travail pour y échapper, mais un salarié à mi-temps au SMIC y
laisserait plus du tiers de son salaire brut. On retient donc une seule contrainte étrangère à
l'efficacité : les bas salaires ne doivent pas payer sensiblement plus qu'aujourd'hui. Le reste découle
des chiffres.

Favoriser le travail renvoie à deux décisions : prendre un emploi, qui dépend de l'impôt sur les
premiers euros, et gagner davantage, qui dépend du taux marginal. Pour les familles, le quotient
familial ne sert plus à rien avec un taux unique, puisqu'il n'agit qu'en faisant changer de tranche.
L'aide la moins coûteuse est un montant fixe par enfant : il ne change le taux marginal de personne, et
seul son financement coûte.

\section*{Le système proposé}

Une seule assiette, le revenu après cotisations sociales, sans abattement ni niche : l'abattement de
10~\% sur les salaires et les pensions, le quotient familial, les réductions et crédits d'impôt et la
contribution exceptionnelle sur les hauts revenus disparaissent dans un impôt unique prélevé à la
source (les revenus du capital au prélèvement forfaitaire unique restent hors champ). Trois paramètres le définissent, le premier étant calé pour que la recette reste de 220~Md€ :
\begin{itemize}
\item 21,4~\% sur les revenus d'activité (salaires, bénéfices des indépendants, rémunérations des
gérants) et les revenus fonciers, au-delà d'une franchise de 10~000~€ par adulte ;
\item 26,4~\% sur les pensions de retraite et les allocations chômage, dès le premier euro ;
\item un crédit de 1~500~€ par enfant, versé même à qui ne paie pas d'impôt (27~Md€ pour 18~millions
d'enfants à charge).
\end{itemize}

La franchise est ce qui sépare ce système de l'impôt par tête. Sans elle (variante A du
tableau~\ref{tab:eff}), un taux de 15,2~\% suffit et la perte sèche tombe à 3~Md€, mais un salarié seul
au SMIC paie 850~€ de plus par an. Chaque millier d'euros de franchise oblige à relever le taux
d'environ 0,6~point ; à 10~000~€, le salarié au SMIC paie à peu près ce qu'il paie aujourd'hui (100~€ de
moins). La franchise est individuelle : dans un couple, le conjoint qui reprend un emploi ne paie rien
sur ses 10~000 premiers euros, alors qu'aujourd'hui ce premier euro supporte la CSG et le taux marginal
d'IR du foyer. Elle ne s'applique ni aux pensions ni au chômage : c'est une aide au travail, pas au
revenu.

L'écart de 5~points entre pensions et travail tient à ce que les retraités d'aujourd'hui ne peuvent
plus modifier leur pension : la taxer ne décourage rien. Il pousse aussi vers l'emploi : à 63~ans,
continuer à travailler plutôt que liquider sa retraite rapporte davantage quand la pension est plus
taxée que le salaire, et l'allocation chômage nette baisse par rapport au salaire qu'on retrouverait.
Le modèle, où pensions et allocations ne réagissent pas, pousserait l'écart bien plus loin : les
5~points sont un choix, pas un résultat.

\begin{table}[!htbp]\centering\scriptsize\setlength{\tabcolsep}{3pt}
\caption{Trois systèmes efficaces comparés à aujourd'hui, à recette égale (220~Md€)}\label{tab:eff}
\begin{tabular}{@{}lrrrr@{}}
\toprule
& Aujourd'hui & A & \textbf{B (proposé)} & C \\
\midrule
Taux sur le travail & barème actuel & 15,2~\% & \textbf{21,4~\%} & 24,4~\% \\
Taux sur les pensions et le chômage & barème actuel & 20,2~\% & \textbf{26,4~\%} & 24,4~\% \\
Franchise par adulte & --- & aucune & \textbf{10~000~€ (travail)} & 10~000~€ (tout revenu) \\
\midrule
Perte sèche, Md€ par an (élasticités basses à hautes) & 18 (8 à 31) & 3 (1 à 5) & \textbf{6 (2 à 10)} & 9 (3 à 15) \\
Taux marginal, salarié seul au SMIC & 9,5~\% & 13,5~\% & \textbf{19,0~\%} & 21,7~\% \\
Taux marginal, salarié seul à 50~000~€ brut & 31,9~\% & 13,6~\% & \textbf{19,2~\%} & 21,9~\% \\
Effet redistributif (aujourd'hui = 100) & 100 & 2 & \textbf{55} & 93 \\
\midrule
\multicolumn{5}{@{}l}{\emph{Impôt aujourd'hui, puis variation par an (+ : paie plus)}} \\
Salarié seul au SMIC (21~500~€ brut) & 2~050~€ & +850~€ & \textbf{$-$110~€} & +170~€ \\
Salarié seul, 50~000~€ brut & 9~000~€ & $-$2~270~€ & \textbf{$-$1~650~€} & $-$610~€ \\
Couple de salariés sans enfant, 50~000~€ & 5~520~€ & +1~220~€ & \textbf{$-$300~€} & +430~€ \\
Couple de salariés, deux enfants, 50~000~€ & 4~770~€ & $-$1~030~€ & \textbf{$-$2~550~€} & $-$1~810~€ \\
Retraité seul, 17~000~€ de pension & 0~€ & +3~430~€ & \textbf{+4~490~€} & +1~710~€ \\
Ensemble des pensions (Md€) & 41 & +19,1 & \textbf{+15,5} & +12,9 \\
Dixième décile de foyers (Md€) & --- & $-$30,5 & \textbf{$-$15,0} & $-$7,1 \\
\bottomrule
\end{tabular}
\par\vspace{3pt}\parbox{\linewidth}{\raggedright\scriptsize\color{encre}
Source : DGFiP, IRCOM et déclarations 2042, revenus 2024 ; OpenFisca-France 176.1 pour les foyers
types ; calculs de l'auteur. Note : IR + CSG + CRDS, avant réactions des contribuables. Tous les
systèmes versent 1~500~€ par enfant. Couple : revenus partagés 60/40. Effet redistributif : baisse de
l'indice de Gini due à l'impôt, rapportée à celle d'aujourd'hui.}
\end{table}

\section*{Ce que cela donne}

La perte sèche tombe de 18 à 6~Md€ par an dans le scénario central, de 8 à 2~Md€ avec les élasticités
basses, de 31 à 10~Md€ avec les hautes. Ce gain vient presque entièrement du haut de la distribution :
11 des 12~Md€ sont gagnés sur les 10~\% de foyers les plus aisés, dont le taux marginal moyen passe à
19~\%, contre 33~\% aujourd'hui entre le 90\textsuperscript{e} et le 99\textsuperscript{e} centile et 43~\% au-dessus. Plus bas, un salarié seul à 50~000~€ brut voit son taux
marginal passer de 32 à 19~\% (figure~\ref{fig:eff}). Le revenu du travail déclaré augmenterait
d'environ 33~Md€ (16 à 57), soit 6~Md€ de recette (3 à 10) non comptés dans le calage : de quoi baisser
le taux d'un demi-point.

Pour les familles, le crédit fait l'essentiel. Un couple de salariés avec deux enfants gagne 2~550~€
par an à 50~000~€ de revenu ; à 30~000~€, il reçoit 1~590~€ au lieu de payer 2~860~€. Au même revenu de
50~000~€, un couple sans enfant ne gagne que 300~€ et un salarié seul 1~650~€.

La baisse sur les salaires (12~Md€) et sur les revenus des indépendants et les loyers (6,5~Md€) est
payée par les retraités (15,5~Md€) et les chômeurs (3~Md€). Un retraité seul à 17~000~€ de pension, qui
ne paie aujourd'hui ni IR ni CSG, paierait 4~490~€, soit 26~\% de sa pension. En moyenne par décile, ces
hausses l'emportent : les déciles 2 à 9 paient 200 à 800~€ de plus par foyer ; le dixième reçoit
15~Md€, 3~500~€ par foyer.

\figbloc{ircsg_e01_systeme}{Taux moyen et taux marginal, aujourd'hui et avec le système proposé}{fig:eff}

\section*{Ce qui peut changer la conclusion}

\textbf{L'ampleur du gain repose sur la réaction des hauts revenus.} Ce sont eux qui portent les
élasticités les plus fortes du modèle, de 0,30 à 0,40 dans le scénario central. Or, en haut de la
distribution, ce qui réagit au taux est surtout le calendrier et la forme des revenus déclarés, plus
que le travail lui-même (Saez, Slemrod et Giertz, 2012). Une assiette sans niche ferme une partie de
ces échappatoires, ce qui réduit le gain à attendre d'une baisse de leur taux. Avec les élasticités
basses, le système rapporte 6~Md€ d'efficacité au lieu de 12, alors que le transfert vers le dixième
décile reste de 15~Md€ : il ne dépend pas des élasticités.

\textbf{La surtaxe des pensions n'est sans coût que pour les retraités actuels.} Pour les actifs, une
pension plus taxée diminue ce que rapportent les cotisations ; pour ceux qui l'anticipent, une partie de
la surtaxe redevient un impôt sur le travail. Surtout, faire payer 4~500~€ de plus à un retraité à
17~000~€ de pension est un choc difficile à imposer, même par étapes. La variante C (24,4~\% pour
tous les revenus, franchise de 10~000~€ y compris sur les pensions) le ramène à 1~700~€, au prix de
3~points de taux en plus sur le travail et de 2,6~Md€ de perte sèche en plus. C'est le coût de
l'égalité de traitement entre retraités et actifs.

\textbf{En bas, la franchise crée une marche.} Au-delà de la franchise, le taux marginal saute de 0 à 19~\% du
salaire brut : un salarié à temps plein au SMIC paie 19~\% sur chaque euro de plus au lieu de 9,5~\%
aujourd'hui, en plus du retrait de la prime d'activité (non modélisée, inchangée). Le système
favorise la prise d'un emploi et la progression des revenus moyens et élevés, pas le passage d'un SMIC à
un SMIC et demi. Adoucir cette marche demanderait un taux d'entrée réduit, donc un second taux.

{\footnotesize\color{encre} \emph{Méthode.} Foyers représentatifs par tranche de revenu fiscal de
référence ; perte sèche $\tfrac12\,e\,z\,\tau^2/(1-\tau)$ sur les revenus du travail, marge intensive
seulement, élasticités de 0,15 à 0,40 selon le revenu. Résultats fragiles au-delà d'un million
d'euros.\par}

\end{document}
